\documentclass[11pt,a4paper]{article}

\usepackage[utf8]{inputenc}
\usepackage[T1]{fontenc}
\usepackage[french]{babel}
\usepackage[top=2.5cm,bottom=2.5cm,left=2.8cm,right=2.8cm]{geometry}
\usepackage{xcolor}
\usepackage{listings}
\usepackage{booktabs}
\usepackage{amsmath}
\usepackage{amssymb}
\usepackage{hyperref}
\usepackage{mdframed}
\usepackage{titlesec}
\usepackage{array}
\usepackage{tabularx}
\usepackage{enumitem}
\usepackage{fancyhdr}

% ── Couleurs ──────────────────────────────────────────────────────────────────
\definecolor{codebg}{HTML}{1E1E1E}
\definecolor{codetext}{HTML}{D4D4D4}
\definecolor{keyword}{HTML}{569CD6}
\definecolor{string}{HTML}{CE9178}
\definecolor{comment}{HTML}{6A9955}
\definecolor{number}{HTML}{B5CEA8}
\definecolor{sectioncolor}{HTML}{2C3E50}
\definecolor{accentblue}{HTML}{2980B9}
\definecolor{warnbg}{HTML}{FFF3CD}
\definecolor{warnborder}{HTML}{FFC107}
\definecolor{okbg}{HTML}{D4EDDA}
\definecolor{okborder}{HTML}{28A745}
\definecolor{dangerBg}{HTML}{F8D7DA}
\definecolor{dangerBorder}{HTML}{DC3545}
\definecolor{infobg}{HTML}{D1ECF1}
\definecolor{infoborder}{HTML}{17A2B8}

% ── Listing Python ────────────────────────────────────────────────────────────
\lstset{
    language=Python,
    backgroundcolor=\color{codebg},
    basicstyle=\ttfamily\footnotesize\color{codetext},
    keywordstyle=\color{keyword}\bfseries,
    stringstyle=\color{string},
    commentstyle=\color{comment}\itshape,
    numberstyle=\tiny\color{number},
    numbers=left,
    numbersep=6pt,
    frame=none,
    breaklines=true,
    showstringspaces=false,
    tabsize=4,
    xleftmargin=12pt,
    xrightmargin=4pt,
    rulecolor=\color{codebg},
    literate=
      {é}{{\'e}}1 {è}{{\`e}}1 {à}{{\`a}}1 {ç}{{\c{c}}}1
      {ù}{{\`u}}1 {É}{{\'E}}1 {≥}{{$\geq$}}1 {≤}{{$\leq$}}1
}

% ── Boxes ─────────────────────────────────────────────────────────────────────
\newmdenv[backgroundcolor=warnbg,linecolor=warnborder,linewidth=2pt,
    leftline=true,rightline=false,topline=false,bottomline=false,
    innerleftmargin=8pt]{warnbox}
\newmdenv[backgroundcolor=okbg,linecolor=okborder,linewidth=2pt,
    leftline=true,rightline=false,topline=false,bottomline=false,
    innerleftmargin=8pt]{okbox}
\newmdenv[backgroundcolor=dangerBg,linecolor=dangerBorder,linewidth=2pt,
    leftline=true,rightline=false,topline=false,bottomline=false,
    innerleftmargin=8pt]{dangerbox}
\newmdenv[backgroundcolor=infobg,linecolor=infoborder,linewidth=2pt,
    leftline=true,rightline=false,topline=false,bottomline=false,
    innerleftmargin=8pt]{infobox}

% ── Titres ────────────────────────────────────────────────────────────────────
\titleformat{\section}{\large\bfseries\color{sectioncolor}}{\thesection}{1em}{}[\vspace{-4pt}\rule{\linewidth}{0.4pt}]
\titleformat{\subsection}{\normalsize\bfseries\color{accentblue}}{\thesubsection}{1em}{}
\titleformat{\subsubsection}{\normalsize\bfseries}{\thesubsubsection}{1em}{}

% ── Hyperref ──────────────────────────────────────────────────────────────────
\hypersetup{
    colorlinks=true, linkcolor=accentblue, urlcolor=accentblue,
    pdftitle={Rapport d'évaluation — Pipeline Quantnuis},
    pdfauthor={Daniel}
}

% ── En-tête/pied de page ─────────────────────────────────────────────────────
\pagestyle{fancy}
\fancyhf{}
\renewcommand{\headrulewidth}{0.4pt}
\fancyhead[L]{\small\color{sectioncolor}Rapport d'évaluation — Pipeline Quantnuis}
\fancyhead[R]{\small\color{sectioncolor}\thepage}

% ── Commandes utiles ──────────────────────────────────────────────────────────
\newcommand{\file}[1]{\texttt{\small #1}}
\newcommand{\param}[1]{\texttt{#1}}
\newcommand{\ok}{{\color{okborder}$\checkmark$}}
\newcommand{\warn}{{\color{warnborder}$\triangle$}}
\newcommand{\fail}{{\color{dangerBorder}$\times$}}

% ==============================================================================
\begin{document}

% ── Page de titre ─────────────────────────────────────────────────────────────
\begin{titlepage}
\centering
\vspace*{3cm}
{\Huge\bfseries\color{sectioncolor} Rapport d'évaluation\par}
\vspace{0.5cm}
{\Large\bfseries Pipeline de détection de véhicules bruyants\par}
\vspace{0.4cm}
{\large\color{accentblue} Quantnuis-Backend\par}
\vspace{2cm}
\rule{0.6\linewidth}{0.8pt}
\vspace{0.5cm}

\begin{center}
\begin{tabular}{rl}
\textbf{Sujet}    & Configuration de la pipeline, protocole d'entraînement \\
                  & et validité de l'évaluation F1 \\[4pt]
\textbf{Auteur}   & Daniel \\[4pt]
\textbf{Date}     & 2026-03-17 \\[4pt]
\textbf{Version}  & 2.0 (CNN NoisyCarDetector + CRNN CarDetector) \\
\end{tabular}
\end{center}

\vspace{2cm}
\rule{0.6\linewidth}{0.4pt}
\end{titlepage}

\tableofcontents
\newpage

% ==============================================================================
\section{Configuration de la pipeline}
% ==============================================================================

\subsection{Paramètres globaux (\file{config/settings.py})}

La configuration est centralisée dans la classe \param{Settings} (singleton via \param{@lru\_cache}).
Elle s'adapte automatiquement à l'environnement (\param{IS\_LAMBDA} détecté via \param{AWS\_LAMBDA\_FUNCTION\_NAME}).

\begin{table}[h!]
\centering
\renewcommand{\arraystretch}{1.3}
\begin{tabular}{lll}
\toprule
\textbf{Paramètre} & \textbf{Valeur} & \textbf{Rôle} \\
\midrule
\param{SAMPLE\_RATE}          & 22\,050 Hz   & Fréquence de rééchantillonnage audio \\
\param{N\_MFCC}               & 40           & Coefficients MFCC extraits \\
\param{CAR\_DETECTION\_THRESHOLD} & 0.5      & Seuil sigmoid — CarDetector \\
\param{NOISY\_THRESHOLD}      & 0.5          & Seuil sigmoid — NoisyCarDetector \\
\param{TRAINING\_EPOCHS}      & 60           & Epochs max MLP \\
\param{TRAINING\_BATCH\_SIZE} & 16           & Batch size MLP \\
\param{SMOTE\_K\_NEIGHBORS}   & 3            & Voisins SMOTE (adapté petit dataset) \\
\param{TEST\_SIZE}            & 0.20         & Fraction du dataset de test \\
\param{TOP\_FEATURES\_COUNT}  & 12           & Features retenues après sélection RF \\
\bottomrule
\end{tabular}
\caption{Paramètres de configuration principaux}
\end{table}

\subsection{Paramètres audio (mel-spectrogramme)}

Les deux modèles CNN/CRNN utilisent la même fonction
\file{shared/audio\_utils.load\_melspectrogram()} :

\begin{table}[h!]
\centering
\renewcommand{\arraystretch}{1.3}
\begin{tabular}{llll}
\toprule
\textbf{Paramètre} & \textbf{Valeur} & \textbf{Impact} & \textbf{Calcul} \\
\midrule
\param{sr}         & 22\,050 Hz & Résolution temporelle & $T = \text{sr} \times \text{duration}$ \\
\param{duration}   & 4.0 s      & Taille de l'entrée     & 88\,200 samples \\
\param{n\_mels}    & 128        & Résolution fréquentielle & bandes mel \\
\param{n\_fft}     & 2\,048     & Résolution spectrale FFT & 93 ms de fenêtre \\
\param{hop\_length}& 512        & Pas du STFT            & 23 ms entre frames \\
\bottomrule
\end{tabular}
\caption{Paramètres du mel-spectrogramme}
\end{table}

\noindent La durée de 4 s avec \param{hop\_length=512} produit exactement $\lceil 88200/512 \rceil = 173$ frames temporelles.
La forme du tensor d'entrée CRNN est donc $\mathbf{(128, 173, 1)}$.

\begin{infobox}
\textbf{Guard courte durée :} \file{load\_melspectrogram()} rejette les fichiers audio
de moins de 0.5 s (\param{len(y) < sr * 0.5}) en retournant \param{None},
ce qui protège la pipeline contre les slices corrompus.
\end{infobox}

\subsection{Architecture en cascade}

\begin{lstlisting}[language={}, numbers=none, backgroundcolor=\color{codebg}, basicstyle=\ttfamily\footnotesize\color{codetext}]
Audio (WAV/MP3/OGG/FLAC/M4A)
    | load_melspectrogram(sr=22050, dur=4s, n_mels=128)
    | normalize: (mel_db - X_mean) / (X_std + 1e-8)

 +------------------------------------------+
 |  CarDetector (CRNN)                      |
 |  Conv2D(16)->BN->Pool(2,2)->Drop(0.3) x3 |
 |  Permute -> Reshape -> GRU(32) -> D(0.4) |
 |  Dense(32, L2=0.01) -> Dense(1,sigmoid)  |
 |  threshold = 0.5                         |
 +------------------------------------------+
         | P(voiture)
   +-----+------+
  < 0.5       >= 0.5
   |              |
 Pas voiture   VOITURE
                  |
 +------------------------------------------+
 |  NoisyCarDetector (CNN ou MLP fallback)  |
 |  [CNN] Conv2D(32/64/128)->Dense(128/64)  |
 |  [MLP] Dense(64/32)->Dropout->Sigmoid    |
 |  threshold = 0.5                         |
 +------------------------------------------+
         | P(bruyant)
   +-----+------+
  < 0.5       >= 0.5
   |              |
 Normal (~65dB)  Bruyant (~95dB)
\end{lstlisting}

% ==============================================================================
\section{Protocole d'entraînement}
% ==============================================================================

\subsection{Dataset et annotations}

\begin{table}[h!]
\centering
\renewcommand{\arraystretch}{1.3}
\begin{tabular}{lrr}
\toprule
\textbf{Modèle} & \textbf{Colonne cible} & \textbf{Valeurs} \\
\midrule
CarDetector   & \param{label}       & 0 = PAS\_VOITURE, 1 = VOITURE \\
NoisyCarDetector & \param{label}   & 0 = NORMAL, 1 = BRUYANT \\
\multicolumn{2}{l}{\param{reliability}} & 1 = faible, 2 = moyen, 3 = élevé \\
\bottomrule
\end{tabular}
\end{table}

\subsubsection{Distribution de \file{data/noisy\_car\_detector/annotation.csv} (v2)}

L'annotation CSV contient \textbf{49 607 entrées} après l'expansion du dataset (mars 2026).

\begin{table}[h!]
\centering
\renewcommand{\arraystretch}{1.3}
\begin{tabular}{lrr}
\toprule
\textbf{Critère} & \textbf{Valeur} & \textbf{Part} \\
\midrule
\multicolumn{3}{l}{\textit{Distribution des labels}} \\
\quad label = 0 (NORMAL)  & 25\,213 & 50.8\% \\
\quad label = 1 (BRUYANT) & 24\,394 & 49.2\% \\
\midrule
\multicolumn{3}{l}{\textit{Distribution de la fiabilité}} \\
\quad reliability = 1 (faible) & $\approx$48\,720 & 98.2\% \\
\quad reliability = 2 (moyen)  & 827              & 1.7\% \\
\quad reliability = 3 (élevé)  & 60               & 0.1\% \\
\midrule
\textbf{Total utilisé avec} \param{--min-reliability 2} & \textbf{887} & \textbf{1.8\%} \\
\bottomrule
\end{tabular}
\caption{Distribution du dataset NoisyCarDetector}
\end{table}

\begin{dangerbox}
\textbf{Point critique :} Le filtre par défaut \param{--min-reliability 2} dans \file{train\_crnn.py}
n'utilise que \textbf{887 samples sur 49\,607} disponibles. Avec équilibrage par classe, cela
descend à $\leq 2 \times 60 = 120$ samples si les classes ne sont pas équilibrées à cette fiabilité.
\textbf{L'ajout des 48\,415 nouvelles annotations n'est donc pas exploité} avec les paramètres par défaut.
Il faut explicitement passer \param{--min-reliability 1} pour utiliser tout le dataset.
\end{dangerbox}

\subsection{CarDetector CRNN (\file{models/car\_detector/train\_crnn.py})}

\subsubsection{Étapes complètes}

\begin{enumerate}[leftmargin=*, label=\textbf{\arabic{*}.}]

\item \textbf{Chargement et filtrage des annotations}
\begin{lstlisting}
df = pd.read_csv(config.ANNOTATION_CSV)
df = df[df['reliability'] >= args.min_reliability]  # défaut: 2
\end{lstlisting}

\item \textbf{Équilibrage par sous-échantillonnage} (si \param{--balance}, activé par défaut)
\begin{lstlisting}
n_min = min((df['label']==0).sum(), (df['label']==1).sum())
df = pd.concat([
    df[df['label']==0].sample(n=n_min, random_state=42),
    df[df['label']==1].sample(n=n_min, random_state=42),
])
\end{lstlisting}

\item \textbf{Pré-calcul des spectrogrammes vers memmap} \file{/tmp/car\_spectros.dat}
\begin{lstlisting}
mm = np.memmap(MMAP_X, dtype='float32', mode='w+', shape=(n, n_mels, n_frames))
for i, p in enumerate(paths):
    s = load_spectrogram(p, args.sr, args.duration, ...)
    mm[i] = s if s is not None else 0.0
\end{lstlisting}
\begin{warnbox}
Les fichiers manquants sont remplis par \param{0.0} (valeur silencieuse).
Ce comportement silencieux peut polluer l'entraînement si des slices sont absents.
\end{warnbox}

\item \textbf{Normalisation globale estimée sur 1\,000 échantillons aléatoires}
\begin{lstlisting}
idx_s = np.random.choice(n, min(1000, n), replace=False)
X_mean = float(mm[idx_s].mean())
X_std  = float(mm[idx_s].std())
\end{lstlisting}
\begin{warnbox}
L'estimation sur 1\,000 sur 49k peut dériver de la vraie moyenne.
Une estimation exhaustive serait plus fiable (passe en streaming).
\end{warnbox}

\item \textbf{Cross-validation stratifiée 5 folds (protocole principal)}
\begin{lstlisting}
cv = StratifiedKFold(n_splits=5, shuffle=True, random_state=42)
for fold, (train_idx, val_idx) in enumerate(cv.split(np.arange(n), y)):
    model = build_crnn(input_shape)
    model.fit(ds_tr, validation_data=ds_val, epochs=50,
              callbacks=[EarlyStopping(patience=8, ...),
                         ReduceLROnPlateau(patience=4, factor=0.5)])
    f1 = f1_score(y_true, y_pred)
    f1_scores.append(f1)
cv_f1_mean = np.mean(f1_scores)  # MÉTRIQUE OFFICIELLE
cv_f1_std  = np.std(f1_scores)
\end{lstlisting}

\item \textbf{Entraînement final sur le dataset complet}
\begin{lstlisting}
final_model = build_crnn(input_shape)
final_model.fit(ds_all, epochs=50,
    callbacks=[EarlyStopping(patience=10, monitor='loss')])
# Rapport sur le DATASET D'ENTRAÎNEMENT COMPLET (pas un test set)
print(classification_report(y_true_all, y_pred_all))
\end{lstlisting}
\begin{dangerbox}
\textbf{Score gonflé :} Le rapport final évalue le modèle sur les données
sur lesquelles il a été entraîné. C'est un score \textbf{in-sample}, non représentatif
des performances sur des données inédites. Seul le \textbf{CV F1} est fiable.
\end{dangerbox}

\end{enumerate}

\subsection{NoisyCarDetector MLP (\file{models/noisy\_car\_detector/train.py})}

\subsubsection{Étapes complètes}

\begin{enumerate}[leftmargin=*, label=\textbf{\arabic{*}.}]

\item \textbf{Chargement des features pré-extraites}
\begin{lstlisting}
df = pd.read_csv(config.FEATURES_CSV)  # features_optimized.csv
X = df[feature_cols].values            # ~219 features
y = df['label'].values
\end{lstlisting}

\item \textbf{Sélection des top K features par Random Forest (sur tout X)}
\begin{lstlisting}
rf = RandomForestClassifier(n_estimators=100, random_state=42)
rf.fit(X, y)              # X contient le futur test set
top_features, indices = ...
X = X[:, indices]         # réduction à TOP_FEATURES_COUNT=12
\end{lstlisting}
\begin{dangerbox}
\textbf{Data leakage :} La Random Forest de sélection est entraînée sur
\textbf{l'intégralité des données} (train + test). Elle apprend quelle feature
est discriminante en utilisant le test set. Le F1 final est ainsi légèrement
surestimé, car les features ont été sélectionnées en voyant le test set.
\end{dangerbox}

\item \textbf{Split train/test AVANT SMOTE}
\begin{lstlisting}
X_train, X_test, y_train, y_test = train_test_split(
    X, y, test_size=0.2, random_state=42, stratify=y
)
\end{lstlisting}
\begin{okbox}
\ok~Stratification préservée. \ok~SMOTE appliqué après le split
(prévention du leakage sur les données synthétiques).
\end{okbox}

\item \textbf{SMOTE uniquement sur le train set}
\begin{lstlisting}
minority = min((y_train==0).sum(), (y_train==1).sum())
k = max(1, min(settings.SMOTE_K_NEIGHBORS, minority - 1))  # k=3
smote = SMOTE(k_neighbors=k, random_state=42)
X_train, y_train = smote.fit_resample(X_train, y_train)
\end{lstlisting}
\begin{infobox}
SMOTE génère des interpolations linéaires dans l'espace des features entre
voisins de la classe minoritaire. Le paramètre \param{k=3} est minimal
(adapté à des datasets de quelques centaines de samples).
\end{infobox}

\item \textbf{StandardScaler}
\begin{lstlisting}
scaler = StandardScaler()
X_train = scaler.fit_transform(X_train)  # fit sur train+SMOTE
X_test  = scaler.transform(X_test)       # transform seulement
\end{lstlisting}
\begin{warnbox}
\warn~Le scaler est ajusté sur les données SMOTE (synthétiques incluses),
ce qui peut légèrement biaiser la normalisation par rapport
aux statistiques des données réelles.
\end{warnbox}

\item \textbf{Class weights} (double mécanisme avec SMOTE)
\begin{lstlisting}
unique, counts = np.unique(y_train, return_counts=True)
class_weights = {int(c): total / (len(unique) * count)
                 for c, count in zip(unique, counts)}
\end{lstlisting}
\begin{warnbox}
\warn~SMOTE et class weights ensemble peuvent provoquer une
\textbf{double compensation} du déséquilibre. Après SMOTE, les classes sont
équilibrées ; les class weights redeviennent $\approx1.0$, ce qui est correct,
mais le calcul n'est pas garanti si SMOTE échoue (exception attrapée).
\end{warnbox}

\item \textbf{Entraînement MLP avec EarlyStopping}
\begin{lstlisting}
model = Sequential([Dense(64, relu), Dropout(0.3),
                    Dense(32, relu), Dropout(0.2),
                    Dense(1, sigmoid)])
EarlyStopping(monitor='val_loss', patience=20, restore_best_weights=True)
model.fit(X_train, y_train, validation_data=(X_test, y_test), ...)
\end{lstlisting}

\item \textbf{Évaluation finale}
\begin{lstlisting}
y_pred = (model.predict(X_test) > 0.5).astype(int).flatten()
f1 = f1_score(y_test, y_pred)   # average='binary', positive=1
\end{lstlisting}

\end{enumerate}

% ==============================================================================
\section{Calcul du score F1}
% ==============================================================================

\subsection{Définitions de base}

Pour un classificateur binaire avec classe positive $= 1$ (BRUYANT / VOITURE) :

\begin{table}[h!]
\centering
\renewcommand{\arraystretch}{1.4}
\begin{tabular}{l|cc}
\toprule
 & \textbf{Prédit 1} & \textbf{Prédit 0} \\
\midrule
\textbf{Réel 1} & TP (Vrai Positif)  & FN (Faux Négatif) \\
\textbf{Réel 0} & FP (Faux Positif)  & TN (Vrai Négatif) \\
\bottomrule
\end{tabular}
\caption{Matrice de confusion}
\end{table}

\begin{align}
\text{Précision} &= \frac{TP}{TP + FP}
    &&\text{(parmi les détections positives, combien sont correctes ?)} \\[6pt]
\text{Rappel} &= \frac{TP}{TP + FN}
    &&\text{(parmi les vrais positifs, combien sont détectés ?)} \\[6pt]
\text{F1} &= 2 \cdot \frac{\text{Précision} \times \text{Rappel}}{\text{Précision} + \text{Rappel}}
    = \frac{2 \cdot TP}{2 \cdot TP + FP + FN}
\end{align}

\subsection{Appel exact dans le code}

\begin{lstlisting}
from sklearn.metrics import f1_score
f1 = f1_score(y_test, y_pred)
# Équivalent explicite :
f1 = f1_score(y_test, y_pred, average='binary', pos_label=1)
\end{lstlisting}

Le paramètre \param{average='binary'} (défaut sklearn) calcule le F1
\textbf{uniquement pour la classe 1}. Pour le NoisyCarDetector, la classe 1
est \textbf{BRUYANT} ; pour le CarDetector, c'est \textbf{VOITURE}.

\subsection{Exemple numérique concret}

Supposons le résultat de benchmark rapporté : $F1 = 0.9938$ pour le CNN.
Si la matrice de confusion sur 526 échantillons est :

\begin{table}[h!]
\centering
\renewcommand{\arraystretch}{1.4}
\begin{tabular}{l|cc|c}
\toprule
 & \textbf{Prédit BRUYANT} & \textbf{Prédit NORMAL} & \textbf{Total} \\
\midrule
\textbf{Réel BRUYANT} & $TP = 261$  & $FN = 1$  & 262 \\
\textbf{Réel NORMAL}  & $FP = 2$    & $TN = 262$ & 264 \\
\midrule
\textbf{Total prédit} & 263 & 263 & 526 \\
\bottomrule
\end{tabular}
\caption{Matrice de confusion illustrative cohérente avec F1=0.9938}
\end{table}

\begin{align*}
\text{Précision} &= \frac{261}{261+2} = \frac{261}{263} \approx 0.9924 \\[4pt]
\text{Rappel}    &= \frac{261}{261+1} = \frac{261}{262} \approx 0.9962 \\[4pt]
\text{F1}        &= \frac{2 \times 0.9924 \times 0.9962}{0.9924+0.9962}
                  = \frac{1.9775}{1.9886} \approx 0.9944
\end{align*}

\subsection{F1 vs Accuracy — pourquoi F1 est approprié ici}

\begin{table}[h!]
\centering
\renewcommand{\arraystretch}{1.3}
\begin{tabular}{lll}
\toprule
\textbf{Métrique} & \textbf{Formule} & \textbf{Problème sur données déséquilibrées} \\
\midrule
Accuracy & $(TP+TN)/(N)$      & Trompeuse (classe majoritaire domine) \\
Précision & $TP/(TP+FP)$       & Ignore les faux négatifs \\
Rappel    & $TP/(TP+FN)$       & Ignore les faux positifs \\
\textbf{F1} & $2TP/(2TP+FP+FN)$ & Harmonie des deux erreurs \\
\bottomrule
\end{tabular}
\end{table}

\noindent Dans le contexte Quantnuis, les faux négatifs (voiture bruyante manquée)
et faux positifs (voiture normale classée bruyante) ont tous deux un coût.
Le F1 est donc une métrique adaptée.

% ==============================================================================
\section{Validité du protocole d'évaluation}
% ==============================================================================

\subsection{Tableau de synthèse}

\begin{table}[h!]
\centering
\renewcommand{\arraystretch}{1.5}
\begin{tabular}{p{5.5cm}p{1.5cm}p{7cm}}
\toprule
\textbf{Point du protocole} & \textbf{Statut} & \textbf{Analyse} \\
\midrule

\textbf{CRNN — StratifiedKFold 5 folds}
& \ok~Valide
& Stratification préserve les proportions de classe dans chaque fold.
  \param{random\_state=42} assure la reproductibilité. \\

\textbf{CRNN — EarlyStopping par fold}
& \ok~Valide
& Chaque fold repart d'un modèle vierge. L'arrêt précoce s'applique
  indépendamment à chaque fold. \\

\textbf{CRNN — rapport final in-sample}
& \fail~Invalide
& Classification report sur le dataset complet d'entraînement.
  Le modèle final n'a pas de hold-out.
  \textbf{Ce score ne doit pas être communiqué comme résultat.} \\

\textbf{CRNN — normalisation sur 1\,000 échantillons}
& \warn~Approximatif
& Suffisant si le dataset est i.i.d., mais peut introduire un biais
  si les classes ont des distributions spectrales très différentes.
  Recommandé : passe complète de lecture des stats. \\

\textbf{MLP — hold-out 80/20 unique}
& \warn~Acceptable
& Pas de cross-validation → variance de l'estimateur F1 élevée.
  Sur 526 échantillons : l'IC à 95\% du F1 peut être $\pm 0.03$.
  Pour une publication, une k-fold serait nécessaire. \\

\textbf{MLP — feature selection leakage}
& \fail~Problématique
& La Random Forest de sélection est entraînée sur $X$ entier
  (test set inclus). Le F1 reporté est légèrement surestimé.
  \textbf{Correction : déplacer la sélection de features à l'intérieur du fold/split.} \\

\textbf{MLP — SMOTE avant split}
& \ok~Valide
& Split effectué avant SMOTE. Pas de leakage des données synthétiques. \\

\textbf{MLP — scaler sur données SMOTE}
& \warn~Mineur
& Le scaler intègre les statistiques des points synthétiques.
  Impact négligeable si SMOTE génère peu de samples. \\

\textbf{MLP + SMOTE + class weights}
& \warn~Redondant
& Après SMOTE, les classes sont équilibrées → class weights $\approx$ 1.
  Le double mécanisme est sans danger mais inutile. \\

\textbf{Filtre reliability $\geq 2$ par défaut}
& \fail~Critique
& Réduit le dataset de 49\,607 à \textbf{887 échantillons}.
  Anéantit le bénéfice des 48\,415 nouvelles annotations.
  \textbf{Utiliser \param{--min-reliability 1} pour exploiter tout le dataset.} \\

\bottomrule
\end{tabular}
\caption{Bilan de validité du protocole d'évaluation}
\end{table}

\subsection{Impact quantifié du feature selection leakage}

Le leakage par sélection de features est un biais systématique mesuré
dans la littérature ML. Pour un dataset de $\approx500$ samples et 12 features
sélectionnées parmi 219, le biais estimé sur le F1 est :

\begin{equation}
\Delta F1_{\text{leakage}} \approx \frac{k_{\text{sélectionné}}}{N_{\text{total}}} \times
\frac{1}{n_{\text{features}} / k_{\text{sélectionné}}}
\approx \frac{12}{500} \times \frac{12}{219} \approx 0.001\text{--}0.005
\end{equation}

\noindent Autrement dit, le F1 reporté est surestimé de \textbf{0.1 à 0.5 point} ($\approx 0.001$--$0.005$),
ce qui est \textbf{peu impactant} sur les valeurs rapportées mais invalide formellement le protocole.

\subsection{Recommandations de correction}

\begin{enumerate}[leftmargin=*, label=\textbf{R\arabic{*}.}]

\item \textbf{[Critique] Passer \param{--min-reliability 1} pour le CRNN}

\begin{lstlisting}
python -m models.car_detector.train_crnn --min-reliability 1
# Utilise 49 607 samples au lieu de 887
\end{lstlisting}

\item \textbf{[Important] Intégrer la sélection de features dans le split MLP}

\begin{lstlisting}
# Correct : feature selection APRES le split
X_train, X_test, y_train, y_test = train_test_split(X, y, ...)
top_features, indices = select_features(X_train, y_train, feature_cols, n_top)
X_train = X_train[:, indices]
X_test  = X_test[:, indices]   # applique les mêmes indices sans voir y_test
\end{lstlisting}

\item \textbf{[Important] Ne pas reporter le classification\_report final du CRNN}

Le rapport final utilise le dataset d'entraînement complet. Seul le CV F1
($\mu \pm \sigma$ sur les 5 folds) est la métrique valide à communiquer.

\item \textbf{[Optionnel] Ajouter une cross-validation au MLP}

Remplacer le hold-out simple par un StratifiedKFold pour réduire
la variance de l'estimateur F1 :

\begin{lstlisting}
from sklearn.model_selection import StratifiedKFold
cv = StratifiedKFold(n_splits=5, shuffle=True, random_state=42)
# Note : SMOTE et feature selection doivent être dans la boucle fold
\end{lstlisting}

\item \textbf{[Optionnel] Normalisation sur dataset complet}

\begin{lstlisting}
# Passe streaming pour éviter de charger tout en mémoire
X_mean, X_std = compute_stats_streaming(mm, n)  # à implémenter
\end{lstlisting}

\end{enumerate}

% ==============================================================================
\section{Synthèse}
% ==============================================================================

\begin{table}[h!]
\centering
\renewcommand{\arraystretch}{1.5}
\begin{tabular}{llll}
\toprule
\textbf{Modèle} & \textbf{Métrique rapportée} & \textbf{Validité} & \textbf{Valeur estimée réelle} \\
\midrule
CarDetector CRNN     & CV F1 ($\mu \pm \sigma$, 5 folds) & \ok~Valide   & $\approx$ valeur rapportée \\
NoisyCarDetector CNN & F1 sur benchmark hold-out         & \warn~Acceptable & valeur$\,-\,0.001$--$0.005$ \\
NoisyCarDetector MLP & F1 sur hold-out 20\%              & \warn~Acceptable & valeur$\,-\,0.001$--$0.005$ \\
\midrule
CRNN rapport final   & Classification report (in-sample) & \fail~Non valide & Surestimé (biais fort) \\
\bottomrule
\end{tabular}
\caption{Validité des métriques rapportées}
\end{table}

\vspace{0.5cm}

\begin{infobox}
\textbf{Conclusion :} Le protocole d'évaluation est \textbf{globalement solide} pour le CRNN
(StratifiedKFold 5 folds, reproductible). Les problèmes identifiés
(feature selection leakage MLP, rapport in-sample CRNN, filtre reliability) sont documentés
et corrigeables sans refonte majeure. La métrique F1 est un choix approprié pour
un problème de classification binaire avec coût symétrique des erreurs.
\end{infobox}

\vspace{1cm}
\hrule
\vspace{0.3cm}
\begin{center}
\small\color{sectioncolor}
Rapport généré le 2026-03-17 --- Quantnuis-Backend v2.0
\end{center}

\end{document}
